\section{Problem Formulation}
\label{sec:problem}
We study \emph{state-channel toxicity propagation} in multi-agent LLM systems: how toxic influence persists in evolving interaction state, which channels carry it, and which mitigations close them. Unlike per-message toxicity evaluation or classical parameter-level unlearning, this focuses on stateful behavioral propagation through transcripts, compressed memory, and model register-matching.

\subsection{Agent-Mediated Conversations as State Machines}
\label{subsec:setup_notation}
We model an agent-mediated discussion as a directed graph $G = (V, E)$ generated over a shared, evolving state.
The process begins with a human seed post $s$ and proceeds through a sequence of agent-generated messages.
Each node $v \in V$ represents a message $x_v$ produced by an agent $a_v \in \mathcal{A}$, and each edge $(u \to v) \in E$ indicates that $x_v$ was generated conditioned on $x_u$ .
The authorship map $v \mapsto a_v$ is \emph{many-to-one}: a single agent may author multiple nodes in the graph.
This matters for the focal agent $A_1$ defined below, which under the \emph{multi-injection} condition authors several replies along a thread rather than only the first response.

At each generation step $t$, the acting agent observes a \emph{conditioning set} $\mathcal{C}(v)$ drawn from the current state and produces a message:
\begin{equation}
    x_v \sim \pi_\theta(\cdot \mid \mathcal{C}(v), r_v),
\end{equation}
where $\pi_\theta$ is the agent's policy (parameterized by an LLM with parameters $\theta$), $\mathcal{C}(v)$ is the visible context, and $r_v$ denotes the agent's role.
The state is then updated to include $x_v$, making it available to future agents. This observe--generate--update loop creates a closed-loop dynamical system in which each agent's output becomes part of the conditioning context for subsequent agents.

\subsection{Toxic Influence as Stateful Contamination}
 
We define the \emph{focal agent} $A_1$ as the first responder to the seed post and vary $A_1$'s behavior:
$A_1 \in \{\texttt{neutral},\ \texttt{toxic}\}$,
while keeping all other components fixed, and each message is scored by an open-source toxicity model for reproducibility (Detoxify~\citep{Detoxify}): $\mathrm{tox}(v) = f_{\mathrm{tox}}(x_v) \in [0, 1]$
% \begin{equation}
%     \mathrm{tox}(v) = f_{\mathrm{tox}}(x_v) \in [0, 1].
% \end{equation}
where larger values indicate more toxic content.
A threshold $\tau{=}0.5$ flags messages as overtly toxic.

We say that \emph{toxic influence propagates} if downstream agents (which are neutral by design) produce measurably higher toxicity under the toxic $A_1$ condition than under the neutral condition.
To isolate downstream propagation from $A_1$'s own contribution, we restrict the average to nodes \emph{not authored by $A_1$}.
Let $V_{A_1} = \{v \in V : a_v = A_1\}$ denote the set of focal-agent nodes (a singleton in the single-injection case, larger under multi-injection).
The downstream toxicity of a graph is then
\begin{equation}
    \mu(G) = \frac{1}{|V \setminus V_{A_1}|} \sum_{v \in V \setminus V_{A_1}} \mathrm{tox}(v),
\end{equation}
and the \textbf{paired effect size} for seed $s$ is:
\begin{equation}
    \Delta\mu(s) = \mu\!\left(G^{(\texttt{toxic})}(s)\right) - \mu\!\left(G^{(\texttt{neutral})}(s)\right).
    \label{eq:effect_size}
\end{equation}
This definition guarantees that any positive $\Delta\mu$ reflects a change in non-focal agents' behavior, not the mechanically high toxicity of $A_1$'s own message.
A positive $\Delta\mu$ therefore indicates that $A_1$'s toxicity has causally influenced the downstream agents' generations.

\subsection{Three Channels of Toxic Persistence}
\label{sec:channels}
We identify three distinct channels through which toxic influence persists in agentic systems:

\noindent \textbf{(i) Transcript backflow.}
Toxic content produced by $A_1$ remains in the raw conversation history and re-enters downstream agents' context at every subsequent turn. Even if a downstream agent would not independently produce toxic content, conditioning on a transcript that contains a high-toxicity message can elevate the Detoxify score of its own output. This is the dominant channel in systems where agents see the full conversation history.

\noindent \textbf{(ii) Memory laundering.}
When agents maintain compressed memory representations (e.g., running summaries of the conversation), toxic framing can be \emph{laundered} through summarization.
The resulting summary may not score as explicitly toxic on classifiers, for instance, a summary might read ``\emph{the discussion has become heated, with participants expressing strong disagreement}'' --- yet conditioning on it raises the expected Detoxify score of subsequent generations relative to conditioning on the matched neutral summary.
We formalize this as:
\begin{equation}
    \mathrm{tox}(M_t) < \tau \quad \text{but} \quad \mathbb{E}\!\left[\mathrm{tox}(v) \mid M_t \in \mathcal{C}(v)\right] > \mathbb{E}\!\left[\mathrm{tox}(v) \mid M_t^{(\texttt{neutral})} \in \mathcal{C}(v)\right],
    \label{eq:laundered}
\end{equation}
where $M_t$ is the contaminated memory and $M_t^{(\texttt{neutral})}$ is the memory from the paired neutral run.

\noindent \textbf{(iii) Parametric bias.}
The model can amplify contamination through register-matching: when toxic transcript or laundered memory re-enters the conditioning context, the policy may produce higher-toxicity outputs than it would under matched neutral context, even after safety fine-tuning.

A robust mitigation must therefore prevent relapse under all three channels simultaneously --- a property we formalize as \emph{backflow-robust mitigation} in~\cref{app:formal_definitions}. Standard LLM unlearning methods (parameter editing, gradient ascent on forget sets) achieve $\Delta\mu \approx 0$ under single-turn evaluation but, as standalone mitigations, are not backflow-robust: they intervene only on \emph{parametric bias} (channel iii) and leave \emph{transcript backflow} (channel i) and \emph{memory laundering} (channel ii) open, allowing toxic content to re-enter the state and trigger relapse.

\subsection{Objective}
\label{sec:objective}

Our objective is to suppress downstream toxic propagation after intervention while preserving the agent's conversational utility. We therefore seek a system that drives the post-intervention paired effect toward zero,
$\mathbb{E}_s[\Delta\mu^{(\mathrm{post})}(s)] \approx 0$,
without materially degrading coherence, topical relevance, or response quality. A formal version of this objective is given in \cref{app:formal_objective}.

This objective requires more than one intervention point. Parameter-level adaptation can reduce amplification of toxic context but cannot remove contaminated transcript or memory from external state; state controls can sanitize reads and gate writes, but residual conflict cues may survive. Memory-only controls likewise miss transcript backflow. This motivates a three-pathway mitigation framework: a fine-tuned policy $\pi_{\theta'}$ for residual parametric amplification, a read-side sanitizer $S$ applied before generation, and a write-side gate $W$ applied before content re-enters transcript or memory.